\section{Conclusion}

We present \solaris, a multiplayer video world model capable of generating consistent multi-view observations from coordinated multi-agent interactions. By developing a scalable multiplayer data collection system, a staged training pipeline transitioning from single-player to multiplayer modeling, and \oursf for memory-efficient long-horizon training, we demonstrated that coherent multi-agent world simulation is achievable.

While this paper leveraged \solarisengine to collect two-player Minecraft data, the platform's potential extends well beyond this setting. It naturally supports more than two concurrent players and can serve a variety of downstream research directions: generating multimodal training data for vision-language models or vision-language-action models, training unified models that jointly perceive and act, developing single-agent and multi-agent policies, studying neurosymbolic approaches where agents reason through code, and constructing benchmarks for 3D understanding, planning, and memory.

Several limitations of our current work also point to promising avenues for future research. First, our training data is entirely synthetic, which introduces gaps in both the action and visual distributions the model encounters, limiting its ability to generalize. While we have demonstrated the value of single-player pretraining, future work could further investigate how to leverage the more abundant single-player data to close this gap. Second, our model lacks persistent memory: when players leave each other's field of view, the model loses track of their shared context, and their trajectories begin to diverge. Unlike a true game engine, the world has no underlying persistent state, as it is specified only through two initial frames, with no mechanism to control or maintain the broader environment. Our work is a valuable starting point for addressing these future research challenges. 
